\subsection{Момент импульса}
%%%%%%%%%%%%%%%%%%%%%%%%%%%%%%%%%%%%%%%%%%%%%%%%%%%%
%%%%%%%%%%%%%%%%%%%%%%%%%%%%%%%%%%%%%%%%%%%%%%%%%%%% 
\z[{\begin{wrapfigure}[5]{r}{0.2\textwidth}%
		\vspace{-12pt}%
		\begin{flushright}%
 			\includegraphics[width=0.2\textwidth]{8.5.1}%
		\end{flushright}%
	\end{wrapfigure}%
	}]%
	{Два диска с моментами инерции $J_1$ и $J_2$ вращаются с угловой скоростью соответственно $\omega_1$ и $\omega_2$ вокруг одной и той же оси без трения. Диски пришли в соприкосновение друг с другом. Из-за возникшего между дисками трения через некоторое время проскальзывание одного диска по другому прекращается. Какова станет тогда угловая скорость вращения дисков? Какое количество теплоты $Q$ выделится?}
%
{$\omega = \dfrac{J_1 \omega_1 + J_2 \omega_2}{J_1 + J_2};$ ~ $Q = \dfrac{J_1 J_2 (\omega_2 - \omega_1)^2}{2(J_1 + J_2)}$}
%%%%%%%%%%%%%%%%%%%%%%%%%%%%%%%%%%%%%%%%%%%%%%%%%%%%	
\z[{\begin{wrapfigure}[5]{r}{0.3\textwidth}%
		\vspace{-10pt}%
		\begin{flushright}%
 			\includegraphics[width=0.3\textwidth]{8.5.2}%
		\end{flushright}%
	\end{wrapfigure}%
	}]%
	{В цилиндр массы $m_1$ и радиуса $R$, покоящийся на гладкой горизонтальной плоскости, попадает пуля массы $m_2$, летящая горизонтально на высоте $h$ от оси цилиндра со скоростью $v$. Считая удар абсолютно неупругим и $m_2 \gg m_1$, найдите скорость оси $v_{\text{ос}}$ и угловую скорость $\omega$ цилиндра.}
%
{$u \approx \dfrac{m_2 v}{m_1};$ ~ $\omega = \dfrac{2 m_2 v h}{m_1 R^2}$.}
%%%%%%%%%%%%%%%%%%%%%%%%%%%%%%%%%%%%%%%%%%%%%%%%%%%% 
\z{На краю свободно вращающегося с угловой скоростью $\omega$ вокруг вертикальной оси диска, имеющего радиус $R$ и момент инерции $J$, стоит человек массы $m$. Как изменится угловая скорость вращения диска, если человек перейдет от края диска к центру? Как изменится кинетическая энергия системы $\Delta K$ при этом? Размерами человека по сравнению с размерами диска пренебречь.}
%
{$\Delta \omega = \dfrac{\omega m R^2}{J};$ ~ возрастет в $\left(1 + \dfrac{m R^2}{J}\right)$ раз.}
%%%%%%%%%%%%%%%%%%%%%%%%%%%%%%%%%%%%%%%%%%%%%%%%%%%%	
\z{Вращающийся обруч радиуса $R$ падает вертикально на горизонтальную плоскость и отскакивает от нее со скоростью $v$ под углом $30\degree$, уже не вращаясь. Какова угловая скорость $\omega_0$ обруча до удара?}
%
{$\omega = \dfrac{v}{2R}$ или $\omega = \dfrac{v \sqrt{2}}{3R}$ (если угол не тот считать).}
%%%%%%%%%%%%%%%%%%%%%%%%%%%%%%%%%%%%%%%%%%%%%%%%%%%% 
\z{По гладкой горизонтальной плоскости навстречу друг другу движутся два одинаковых тонких вращающихся кольца. Их скорости $v_1$ и $v_2$ направлены по прямой, соединяющей центры колец. Угловые скорости колец $\omega_1$ и $\omega_2$. Определите угловую скорость $\omega$ колец после соударения, если проскальзывание их относительно друг друга исчезает в последний момент удара.}
%
{$\omega_1' = \dfrac{3\omega_1 - \omega_2}{4};$ ~ $\omega_2' = \dfrac{3\omega_2 - \omega_1}{4}$.}